\section{Method}
\label{sec:method}

\subsection{Two-locus type safety: information flow vs output validity}
\label{sec:method:two_loci}

A type-safe reasoner over an Olog $\mathcal{O}$ has two distinct correctness
properties to enforce, and they live at different layers of the
architecture.

\textbf{(B) Information-flow correctness} (attention layer): the model
should not internally mix information across token pairs whose types are
unrelated under $\mathcal{O}$. This is enforced by gating attention
weights by Olog reachability (\S\ref{sec:method:attention}).

\textbf{(D) Output validity} (sampling boundary): the model's emitted
labels at each generation step should be admissible under $\delta_\mathcal{O}$,
the typed transition relation given by the Olog's morphisms. This is
enforced by a constrained decoder that masks the output distribution at
each step to the admissible-label set (\S\ref{sec:method:decoder}).

The distinction matters: (B) restricts internal representations; (D)
restricts emissions. (B) without (D) is information-flow-clean but does
not guarantee trajectory soundness, because attention reachability admits
the eventual destination type without requiring each step to correspond
to an Olog edge from the current type. (D) without (B) is sound at the
sampling boundary but the internal representations the model learns may
still mix unrelated types. The composition (B)+(D) gives both
information-flow cleanliness and trajectory soundness.

\subsection{(B) Ontology-Typed Attention}
\label{sec:method:attention}

Standard self-attention computes weights $a_{ij} \propto \exp(q_i^\top
k_j)$ over all token pairs. We type each token with an Olog object (or
$\bot$ for untyped tokens) and mask attention by Olog reachability:
\begin{equation}
\tilde{a}_{ij} = a_{ij} \cdot \mathbb{1}[\tau(j) \in
\mathrm{Reach}_{\mathcal{O}}(\tau(i))],
\label{eq:typed_attention}
\end{equation}
where $\tau(\cdot)$ is the Olog type assignment and
$\mathrm{Reach}_{\mathcal{O}}(t)$ is the set of Olog objects reachable from
$t$ under composition of morphisms. Untyped tokens are unrestricted on
either side. Eq.~(\ref{eq:typed_attention}) is the (B) configuration.

\paragraph{What (B) buys, and what it does not.} Reachability is a
transitive closure: from a current type $t$, (B) admits attention to any
type $t'$ such that there exists a composition of Olog morphisms
$t \to \cdots \to t'$. This is the right primitive for representation
hygiene: information flows only along type-respecting paths inside the
model. It is the wrong primitive for output emission: a destination
$t'$ may be many morphisms away from $t$, and the next emitted label
must correspond to an edge $t \to t''$ in the Olog, not just any type
reachable from $t$. Without (D), a (B)-only system can still emit
trajectories that violate $\delta_\mathcal{O}$.

\subsection{(D) Constrained Decoder}
\label{sec:method:decoder}

At each generation step, given the current type $t$ and the proof step
$(t, r, t')$ recommended by the prover, the decoder constructs the
admissible-label set $\mathcal{A}(t, r, t')$ as the union of (i) the
target type $t'$, (ii) the relation token $r$, and (iii) function words
admissible regardless of type (articles, prepositions, punctuation). The
output logits are then masked:
\begin{equation}
\tilde{\ell}_v = \begin{cases}
\ell_v & v \in \mathcal{A}(t, r, t') \\
-\infty & \text{otherwise.}
\end{cases}
\end{equation}
This is the (D) configuration; the implementation is the
\texttt{ConstrainedDecoder} of Appendix~\ref{app:impl}. The decoder
operates on logits, not on hidden states, and is independent of (B);
they compose without interference.

\subsection{Sheaf Cohomology of a KG}
\label{sec:method:sheaf}

We model the KG as a cellular sheaf $\mathcal{F}$ over the underlying
multigraph (\S\ref{sec:background}, full construction in
Appendix~\ref{app:sheaf}). The first cohomology $H^1(\mathcal{F})$ is the
quotient $\ker \delta_1 / \mathrm{im}\,\delta_0$; non-trivial classes
correspond to inconsistency that cannot be resolved by adjusting vertex
values alone. We compute $\dim H^1$ via the sheaf Laplacian and the
ranks of the boundary operators (Appendix~\ref{app:sheaf}).

\begin{theorem}[Conflict Sensitivity, informal]
\label{thm:conflict_sensitivity}
Let $\mathcal{F}$ be a cellular sheaf on a KG and $\mathcal{F}'$ the same
sheaf with $k$ injected conflicts (edges whose restriction maps disagree
on their endpoints). Then $\dim H^1(\mathcal{F}') \geq \dim
H^1(\mathcal{F}) + c$ for some $c \leq k$ depending on conflict
locality.
\end{theorem}
The proof and full hypotheses are deferred to
Appendix~\ref{app:conflict_proof}.

\subsection{Proof-Guided Generation}
\label{sec:method:proof}
Each prediction $(h, r, t)$ is accompanied by a derivation
$\pi: h \xrightarrow{r_1} \cdots \xrightarrow{r_n} t$ in the Olog. The
\emph{prove-then-emit} procedure runs in three steps:
\begin{enumerate}
\item \textbf{Prove.} Search the Olog for a morphism path
$\pi = (r_1, \ldots, r_n)$ from $\tau(h)$ to a goal type compatible
with $r$, by breadth-first search over the morphism set
($O(|T| + |R|)$). If no path exists, abstain.
\item \textbf{Check.} Test whether the candidate triple closes a
non-trivial class in $H^1(\mathcal{F})$ (\S\ref{sec:method:sheaf}); if
so, flag the prediction as structurally inconsistent.
\item \textbf{Emit.} For each step $(t, r_i, t')$ along $\pi$, the (D)
decoder restricts the output logits to the admissible-label set
$\mathcal{A}(t, r_i, t')$ (\S\ref{sec:method:decoder}) and samples the
next label.
\end{enumerate}
The proof object $\pi$ is the shared interface that drives both the
sheaf consistency check (step~2) and the constrained decoder (step~3),
and is also the artifact returned to the user as the derivation behind
the prediction.

\subsection{Pipeline}
\label{sec:method:pipeline}

\begin{figure}[t]
\centering
\begin{tikzpicture}[
  >={Stealth[length=2.2mm]},
  every node/.style={font=\small},
  box/.style={draw, rounded corners, align=center,
              minimum height=11mm, minimum width=21mm, inner sep=3pt},
  stage/.style={box, fill=blue!7},
  io/.style={box, fill=black!6},
  art/.style={box, fill=orange!16},
]
\node[io]    (q)   {Query $(h,r)$};
\node[stage] (b)   [right=7mm of q]  {(B) typed\\attention};
\node[stage] (pv)  [right=7mm of b]  {Olog prover\\(BFS)};
\node[stage] (d)   [right=7mm of pv] {(D) constrained\\decoder};
\node[io]    (out) [right=7mm of d]  {$(h,r,t)$\\$+$ derivation};
\node[art]   (pf)  [below=10mm of pv] {proof object $\pi$};
\node[stage] (sh)  [below=10mm of d]  {sheaf $H^1$\\check};
\draw[->] (q)  -- (b);
\draw[->] (b)  -- (pv);
\draw[->] (pv) -- (d);
\draw[->] (d)  -- (out);
\draw[->] (pv) -- (pf);
\draw[->] (pf) -- (sh);
\draw[->] (pf) -- (d);
\draw[->] (sh) -- (out);
\end{tikzpicture}
\caption{\textbf{The prove-then-emit pipeline.} A query passes through
(B) typed attention (an information-flow restriction inside the model),
an Olog prover that produces the proof object $\pi$, and the (D)
constrained decoder (an output-validity restriction at the sampling
boundary). The proof object is the shared interface: it drives both the
constrained decoder and the sheaf $H^1$ consistency check, which can
flag a structurally inconsistent prediction or route it to an abstain
head. (B) and (D) act at different loci and compose without
interference.}
\label{fig:pipeline}
\end{figure}

Figure~\ref{fig:pipeline} shows the full pipeline; a query $(h, r)$ is
processed in four stages.
(i)~The entity-token sequence is encoded under (B) typed attention, restricting information flow to Olog-reachable type pairs.
(ii)~A prover constructs the derivation $\pi: h \xrightarrow{r_1} \cdots \xrightarrow{r_n} t$, which serves as the proof object for subsequent stages.
(iii)~The sheaf consistency check tests whether the candidate triple $(h, r, t)$ induces a non-trivial cohomology class; predictions that do are flagged as structurally inconsistent or routed to an abstain head.
(iv)~The (D) constrained decoder restricts output logits at each step to the admissible-label set $\mathcal{A}(t, r, t')$ derived from the proof object.
The proof object is the single shared interface between stages (ii)--(iv); it makes the pipeline end-to-end verifiable.
